\subsection{Conformance suite}
\label{sec:conformance-suite}

The conformance suite pairs each test with the verdict given by Python itself, and runs on CPython,
GraalPy, and Pyodide. A test in \filename{python-raises/typing-rejects/}, for example, is a program which raises an exception other than \pyinline{SyntaxError} under Python and is rejected by the \PurePy type checker; the directory of a test names the row of \figref{taxonomy} for the outcome under Python, followed by the column for the outcome under \PurePy. \speconly{A test which fails because of a bug reported in mypy or a Python implementation is recorded as expected to fail; \figref{reported-bugs} lists these behaviours. }\todo{Reference interpreter, needed
to detect dynamically excluded operations (\issue{64}).}

\begin{AddedBlock}
The first line of a test of an ill-formed program cites the form of the innermost judgement to fail in an attempted
derivation, by its identifier in the conformance package, together with the rule applying to the judgement if there
is one, as in \filename{\# judgement: synth, rule: call}. A rule applies if it is the only rule with a conclusion
matching the judgement, so that an unbound variable cites \ruleName{var}, or if it is the only one of several such
rules with its selecting premises holding, such as the type synthesised by a subexpression or the context entry of a
name, so that a class given too many type arguments cites \ruleName{ty-class} and a subscript of a list cites
\ruleName{subscript-list}. The cited rule then has a failing premise or side condition: a call with too few arguments
to a callee of union type cites \ruleName{call}, and \pyinline{1 + "a"} cites \ruleName{binop}, since operator
resolution is a side condition. The judgement form is cited alone if no rule applies: \pyinline{x: Nonexistent = 5}
cites the form for type expressions, and a subscript of an integer cites the form for synthesis. A type mismatch
cites the rule of subtyping matching both types, with a literal type taken as its base type except against a literal
type, so a test assigning a \pyinline{list[str]} to a variable declared at \pyinline{list[int]} cites
\ruleName{subty-list}. If no single rule of subtyping matches both types, the mismatch cites the rule asking for the
check: \pyinline{x: int = "a"} cites \ruleName{declare-assign}, not \ruleName{check-syn}.
\end{AddedBlock}
